\documentclass[11pt]{article}
\usepackage[T1]{fontenc}
\usepackage[margin=1in]{geometry}
\usepackage{amsmath,amssymb,amsthm,microtype}
\usepackage[colorlinks=true,urlcolor=blue,citecolor=blue]{hyperref}
\newtheorem{theorem}{Theorem}
\title{A scalar counterexample to property $(P')$ as printed in arXiv:1003.0108}
\author{Solution packet for \texttt{fa\_banach\_001}}
\date{11 August 2026}

\begin{document}
\maketitle

\noindent\textbf{Status.}
Full negative answer to the literal statement of $(P')$ in Section 6.
The counterexample already lies in the one-point scalar rational setting.

\section*{The printed question}

Ball and Sasane define the stability margin
\[
 \mu_{P,C}=\|H(P,C)\|_\infty^{-1}
\]
when $C$ stabilizes $P$, and define it to be zero otherwise.  Their final
section asks whether the following property always holds in the abstract
setup:
\begin{quote}
$(P')$: $C$ satisfying $\mu_{P_0,C}>m$ stabilizes $P$ only if
$d_\nu(P,P_0)\le m$.
\end{quote}
Read literally, this says
\begin{equation}\label{eq:literal}
 \mu_{P_0,C}>m\quad\hbox{and}\quad C\text{ stabilizes }P
 \quad\Longrightarrow\quad d_\nu(P,P_0)\le m.
\end{equation}

\begin{theorem}
Implication \eqref{eq:literal} is false under assumptions (A1)--(A4), even
with $R=S=\mathbb C$, scalar plants, and $m=\frac12$.
\end{theorem}

\begin{proof}
Take $R=S=\mathbb C$ with complex conjugation as involution, take the index
group to be trivial, and let $\iota=0$.  This satisfies (A1)--(A4):
$\mathbb C$ is a unital integral domain and a semisimple commutative Banach
algebra, and every nonzero element is invertible.

Set
\[
 P_0=0,\qquad P=1,\qquad C=0.
\]
Use the normalized coprime factorizations
\[
 P_0=\frac0{1},\qquad
 P=\frac{1/\sqrt2}{1/\sqrt2},\qquad
 C=\frac0{1}.
\]
Thus
\[
 G_0=\begin{bmatrix}0\\1\end{bmatrix},\qquad
 \widetilde G_0=\begin{bmatrix}-1&0\end{bmatrix},\qquad
 G=\begin{bmatrix}1/\sqrt2\\1/\sqrt2\end{bmatrix}.
\]
The determinant condition in the definition of $d_\nu$ holds because
\[
 G^*G_0=\frac1{\sqrt2}\ne0
\]
and the index is trivial.  Consequently
\[
 d_\nu(P,P_0)=\|\widetilde G_0G\|_\infty
 =\left|{-\frac1{\sqrt2}}\right|=\frac1{\sqrt2}>\frac12.
\]

On the other hand, the source's closed-loop matrix is
\[
 H(Q,C)=\begin{bmatrix}Q\\1\end{bmatrix}(1-CQ)^{-1}
          \begin{bmatrix}-C&1\end{bmatrix}.
\]
For the nominal pair,
\[
 H(P_0,C)=H(0,0)=\begin{bmatrix}0&0\\0&1\end{bmatrix},
 \qquad \|H(0,0)\|=1,
\]
so $C$ stabilizes $P_0$ and $\mu_{P_0,C}=1>\frac12$.  The same
controller also stabilizes $P=1$, since
\[
 H(P,C)=H(1,0)=\begin{bmatrix}0&1\\0&1\end{bmatrix}
 \in R^{2\times2}.
\]
Hence both hypotheses on the left side of \eqref{eq:literal} hold with
$m=\frac12$, while its conclusion fails.  This proves the claim.
\end{proof}

\section*{Interpretation boundary}

The example uses constant transfer functions, so it also belongs to the
rational class that the source says already satisfies $(P')$.  Therefore the
displayed sentence cannot simultaneously have its literal implication
meaning and the asserted rational-case provenance.  Standard robust
stability says that a controller with margin exceeding a radius stabilizes
\emph{every plant in the corresponding $\nu$-ball}; it does not say that a
controller can stabilize no plant outside that ball.  Indeed, the zero
controller stabilizes every stable scalar plant, including many plants far
from the zero plant.

Thus the open question as printed has a complete negative answer.  If the
authors intended a worst-case statement about an entire uncertainty ball, or
a minimality statement among metrics, that is a different quantified problem
and is not refuted by this packet.

\begin{thebibliography}{9}
\bibitem{BS}
J.~A. Ball and A.~J. Sasane,
\emph{Extension of the $\nu$-metric},
Complex Analysis and Operator Theory 6 (2012), 65--89,
arXiv:1003.0108,
\url{https://arxiv.org/abs/1003.0108}.
\end{thebibliography}

\end{document}
